\properlane{conversion-without-rivalry}{It is still time to enter, and the work is mercy}
\subsection{A late invitation asks for a real answer}
Chrysostom identifies the vineyard first with God's commandments and the working time with present life. The hours represent the ages at which people approach God and undertake the good. The controlling purpose is to give courage to those converted in extreme old age: they must not suppose that their late beginning confines them to a lesser portion. This is encouragement toward earnestness. The late workers do not merely receive a consoling message while remaining where they were; they enter the vineyard.\footnote{Chrysostom, \textit{Homilies on Matthew} 64.3--4.}

This reading puts pressure on two very different listeners. The longtime disciple is tempted to make an early beginning a present title to honor. The latecomer is tempted to make a bad past a final verdict against hope. Neither knows the gift rightly while measuring it chiefly by elapsed time. The first must relinquish pride; the second must begin. The call is addressed to persons whose response matters, and the story's generosity makes their response more urgent rather than less necessary.

Augustine joins the same practical conclusion through a different question. If the eleventh hour receives the same denarius, why not deliberately wait for it? Because the promised reward does not promise another hour of life. Those summoned earlier did not announce that they would come later. A person called now cannot convert another person's late mercy into an entitlement to postpone conversion. Augustine then treats despair and presumptuous security as opposite ways of refusing a saving hope: one says forgiveness is impossible, the other that repentance can always be delayed.\footnote{Augustine, Sermon 87, §§7--10.}

Isaiah supplies the verbs this response needs. Seeking God requires forsaking one's former way and thoughts. Jerome connects the opportunity to embodied life, while Aquinas emphasizes seeking before adversity and death. The future gift does not make the present morally empty. It makes the present the place where a return can happen. The Entrance's confidence in help during distress prevents urgency from becoming hopelessness: the hearer is summoned by a Lord who hears, not abandoned to reconstruct life alone.\footnote{Jerome, \textit{Isaiah} XV on 55:6--7; Aquinas, \textit{Isaiah} 55.}

\subsection{The promise heals comparison rather than abolishing work}
Chrysostom's refusal of heavenly envy is central here. He reasons from the saints' charity: those who give their lives for sinners in this life will rejoice when those sinners enjoy salvation. The Gospel complaint therefore dramatizes the greatness of the latecomer's gift. It is not an invitation to imagine resentment as an unavoidable feature of even the holiest community. Present envy can be healed because the promised communion is free of it.\footnote{Chrysostom, \textit{Matthew} 64.3.}

His interpretation also distinguishes the equal wage from the closing first--last maxim. Inside the parable the early workers still receive the agreed payment; the others are unexpectedly made equal. The maxim can point beyond that equality toward an actual reversal, where someone who began well neglects virtue and someone formerly sunk in vice advances. Chrysostom permits a hint concerning Jews in this latter discussion as well as changes among believers. His primary personal-age reading is not an explicit polemic rejecting every Jewish--Gentile application. His received Gospel lemma includes the many-called clause beyond today's endpoint, so the full scope of his argument reaches beyond the end of the shorter appointed text.\footnote{Chrysostom, \textit{Matthew} 64.3--4; today's Gospel is Matthew 20:1--16a.}

The practical consequence is neither complacency nor competition. A person does not imitate Christ by becoming satisfied that someone else is doing worse. Chrysostom's sermon turns from the parable to care for conduct, almsgiving, and the choice of examples. The negligence of a rich neighbor excuses nothing. Christ and those who do good are the examples that can produce humility and diligence. This is a more demanding freedom than the relief of being ahead of somebody: it removes the convenient comparison by which neglect disguises itself as adequacy.\footnote{Chrysostom, \textit{Matthew} 64.5.}

The Collect identifies love of God and neighbor as the work's governing form. Psalm 145 makes mercy toward all God's works an object of praise. Praising that universality while resenting a neighbor's pardon would divide the act of worship against itself. Augustine's exposition insists that God's patience must not become permission to remain in sin. Bellarmine likewise sees patience inviting repentance. These are grounds for a hopeful return and for refusing to obstruct another's return. Mercy is to be received and practiced, not praised abstractly while withheld personally.

\subsection{Attendance becomes hospitality, and desire becomes service}
The adapted acclamation locates conversion at the moment of hearing. Lydia's opened heart attends to the message, and her response becomes hospitality. Chrysostom emphasizes both sides of the event: divine help does not erase her willingness, and her willingness does not erase the Lord's action. The acclamation can therefore be heard as a request for something more searching than concentration. A hearer needs openness to the change that the Word may require.\footnote{Chrysostom, \textit{Acts} 35, on 16:13--15.}

Augustine's account of faith's beginning places a different stress on that same event. The opening for which the Church prays is God's gift; it is not a reward for an independently completed first movement of faith. Their common practical ground is that hearing bears fruit in response. Neither passage warrants treating the assembly as spectators while conversion takes place without them, and neither licenses self-congratulation for a heart allegedly opened without help.\footnote{Augustine, \textit{On the Predestination of the Saints} I.41; Chrysostom, \textit{Matthew} 64.3.}

Paul gives the positive alternative to rivalry. He desires the joy of being with Christ, yet continued life is useful to others. Chrysostom presents the willingness to postpone rest as charity for Christ's sake. A worker formed by that love does not merely refrain from complaining about a latecomer; he labors for the latecomer's good. This changes the moral question from how much advantage one can claim to whom one's life can help. Paul's closing demand for conduct worthy of the Gospel extends the same responsibility to the hearers.

The Eucharistic prayers make that response continuous with worship. In the Prayer over the Offerings, professed faith reaches sacramental participation. In the Prayer after Communion, the life corresponding to redemption remains an object of petition. The gift is real and its reception matters, but receiving does not make conduct irrelevant. Augustine's account of the Lord's table explains why: the recipient of Christ's self-giving becomes responsible for a corresponding charity toward the brethren. The sacramental gift provides a reason and strength to serve, not an exemption from service.\footnote{Augustine, \textit{John} 47.2.}

The Psalm Communion, if used, keeps this response in the grammar of dependence: the person who hears the command asks for steady ways in fulfilling it. Bellarmine explains this petition as a rejection of confidence in unaided power. The John Communion, if used instead, makes response personal: the sheep belong to a Shepherd whose knowledge is joined to self-giving. Conversion then means following a known Lord rather than trying to manufacture virtue in isolation. Both endings prepare the final petition for redemption to bear fruit beyond the moment of reception.

\subsection{Four senses of the timely call}
\begin{fourSenses}
\item[Literal.] The invitations come at different hours, and those called respond. Isaiah's present summons, the psalm's patient mercy, Paul's useful remaining, and Lydia's attention describe real acts in their own contexts. The Entrance promises help, the Collect names charity, and the Eucharistic prayers carry reception toward a lived response. The two Communion alternatives accent either requested obedience or relationship with the Shepherd.

\item[Allegorical.] Christ calls people into the vineyard of his commandments during the working day of life. His word addresses the Church now, and his sacramental gift sustains its charity. The good Shepherd and the caller are not competing masters; the members receiving him are drawn into the pattern of his service.

\item[Moral.] Begin without despair and persevere without entitlement. Augustine's warning against delay guards Chrysostom's encouragement of late converts. Exchange envy for mercy and comparative excuses for concrete almsgiving and service. The opened heart must attend, and the person receiving Communion must live what faith confesses.

\item[Anagogical.] Full hope is offered to the latecomer, and the kingdom's joy contains no envy. Eternal communion releases the hearer from the fear that a late beginning makes genuine conversion pointless. The same hope refuses presumption: today's call is an opportunity to enter, not a guarantee that tomorrow's call will be available.
\end{fourSenses}
